\begin{textsample}{NEG Dim 5 – vem\_ed\_25 – Score 2.00 – vem\_ed\_25/t136.txt}  \label{ex:f5_neg_156}
Artigo de opinião As WebQuests como ferramentas de sala de aula invertida Edilene Gasparini Fernandes Fatec São José do Rio Preto edilene@fatecriopreto.edu.br As WebQuests foram criadas por Bernie Dodge, em 1995, e foram usadas com o objetivo de integrar disciplinas e trazer a web para dentro das salas de aula do ensino fundamental e médio.
Elas podem ser consideradas como uma das ferramentas para a metodologia ativa conhecida como sala de aula invertida (flipped classroom), em que o conteúdo é apresentado fora do ambiente escolar e o tempo em sala de aula é usado pelo professor para tirar dúvidas e promover debates com seus estudantes.
Em 2020, as WebQuests voltaram a ser discutidas no Brasil por meio da pesquisa de Lilian Bacich sobre seu uso dentro do meio acadêmico (Bacich, 2020).
Motivada por esse ressurgimento, propus uma adaptação de seu uso no ensino superior tecnológico das Fatecs com o objetivo de melhoria do ensino-aprendizagem de língua inglesa.
Essa pesquisa foi financiada pelo CPS, como Regime de Jornada Integral, e envolveu alunos de três cursos diferentes da Fatec São José do Rio Preto (Tecnologia em Agronegócio, Tecnologia em Informática para Negócios e Tecnologia em Análise e Desenvolvimento de Sistemas), que colaboraram na criação e hospedagem das WebQuests na plataforma Google Sites.
As WebQuests são ferramentas que provocam o aluno por meio do questionamento, da curiosidade, e o motivam a buscar uma solução para o problema apresentado.
As etapas de uma WebQuest podem variar conforme o alvo de sua questão, mas, classicamente, Bernie Dodge estabelece: 1.
Introdução - fase em que se prepara o “palco” e se “fornece algumas informações de fundo”. 2.
Apresentação de uma tarefa “factível e interessante” que instigará o grupo pela curiosidade e aderência com sua realidade e aplicação. 3.
Indicação de endereços eletrônicos que sejam a base de partida para a aventura que terá outras paragens, naturalmente. 4.
Abordagem da estrutura de etapas na qual será construída a WebQuest: a introdução do problema ou questão, a tarefa dividida entre os membros do grupo, o processo de construção da abordagem e os recursos de que farão uso. continuação Em seguida, o projeto passa por uma \textbf{avaliação} interna para que os estudantes cheguem, juntos, à conclusão.
Depois de concluída, a WebQuest é também avaliada por professores ou responsáveis pelo projeto (Dodge, 1996, p. 1-2).
O resultado positivo que o uso das WebQuests tem proporcionado com os alunos da Fatec São José do Rio Preto ganhou também o ambiente dos Projetos Colaborativos Internacionais (PCIs) realizados na unidade.
Desde 2022, elas fazem parte dos PCIs que a Fatec Rio Preto realiza com países como África do Sul, Índia e Armênia.
De início elas foram empregadas como ferramentas dentro dos projetos colaborativos, mas a proposta de tornar a interação um trabalho colaborativo de criação de novas WebQuests tornou as interações ainda mais animadas.
Professores e alunos superaram-se em criatividade e inovação.
Em outubro de 2024, está prevista mais uma edição do PCI com alunos de Letras da Yerevan State University (Armênia), cuja interação lhes fornece know-how para o trabalho com as Tecnologias da Informação e da Comunicação (TICs) em ambientes colaborativos.
Yan Ribon, aluno da Fatec São José do Rio Preto envolvido com o desenvolvimento das WebQuests, foi agraciado com uma bolsa do CNPq em 2023 a partir de um projeto de um aplicativo para acesso mobile das WebQuests.
O aplicativo WebEnglish acabou de ser finalizado e está disponível, (exclusivamente no sistema operacional Android) aos colegas de trabalho e a todos que se interessarem pelo seu uso, por meio deste link: https://github.com/yanrex8/app-webquests Coloco-me à disposição para orientar quaisquer professores que desejem implementar a ferramenta em suas aulas de línguas, ou ainda, que pretendam trabalhar com elas em PCIs em quaisquer áreas, já que os assuntos que elas abordam são variados e envolvem diferentes áreas de conhecimento.
Referências BACICH, Lilian.
WebQuest: como organizar uma atividade significativa de pesquisa.
Inovação na educação.
São Paulo, 22 de março de 2020.
Disponível em: https://lilianbacich.com/2020/03/22/webquest-como-organizar-uma-atividade-significativa-de-pesquisa/ Acesso em: 12 set. 24 DODGE, Bernie.
Webquest: uma técnica para aprendizagem na rede.
UFSCar, 1996.
Disponível em: https://www.dm.ufscar.br/~jpiton/downloads/artigo\_webquest\_original\_1996\_ptbr.pdf Acesso em: 12 set. 24

% matched lemmas: avaliação
\end{textsample}
